\documentclass[11pt]{article}
% =====================================================================
%  preamble.tex  —  shared preamble for Part I (physics.tex) and
%  Part II (math.tex) of "Bell Correlation as a Foundational Principle".
%
%  RECONSTRUCTED minimal preamble.  The original was referenced in
%  README_build_notes.md but never created as a file; physics/math were
%  only ever compiled against a throwaway stub.  This file collects the
%  macros and environments that physics.tex and math.tex are known to use
%  (from project records): a single shared theorem-numbering scheme across
%  both parts, the status labels, the NSA symbols (st, ord, *R), and the
%  categorical chain C_corr -> C_dist -> C_act -> C_eq.
%
%  BUILD (per README_build_notes.md), bidirectional xr cross-references:
%    each of physics.tex / math.tex does % =====================================================================
%  preamble.tex  —  shared preamble for Part I (physics.tex) and
%  Part II (math.tex) of "Bell Correlation as a Foundational Principle".
%
%  RECONSTRUCTED minimal preamble.  The original was referenced in
%  README_build_notes.md but never created as a file; physics/math were
%  only ever compiled against a throwaway stub.  This file collects the
%  macros and environments that physics.tex and math.tex are known to use
%  (from project records): a single shared theorem-numbering scheme across
%  both parts, the status labels, the NSA symbols (st, ord, *R), and the
%  categorical chain C_corr -> C_dist -> C_act -> C_eq.
%
%  BUILD (per README_build_notes.md), bidirectional xr cross-references:
%    each of physics.tex / math.tex does % =====================================================================
%  preamble.tex  —  shared preamble for Part I (physics.tex) and
%  Part II (math.tex) of "Bell Correlation as a Foundational Principle".
%
%  RECONSTRUCTED minimal preamble.  The original was referenced in
%  README_build_notes.md but never created as a file; physics/math were
%  only ever compiled against a throwaway stub.  This file collects the
%  macros and environments that physics.tex and math.tex are known to use
%  (from project records): a single shared theorem-numbering scheme across
%  both parts, the status labels, the NSA symbols (st, ord, *R), and the
%  categorical chain C_corr -> C_dist -> C_act -> C_eq.
%
%  BUILD (per README_build_notes.md), bidirectional xr cross-references:
%    each of physics.tex / math.tex does \input{preamble} then
%    \usepackage{xr}\externaldocument[I-]{physics}  (in math.tex)
%    \usepackage{xr}\externaldocument[II-]{math}     (in physics.tex)
%    Compile order: bootstrap each once (pdflatex) to make .aux, then
%    latexmk both twice to converge cross-document refs.
%
%  This preamble is deliberately conservative: it defines every macro the
%  two documents are known to reference, plus the common analysis symbols,
%  so a missing \input does not break the build.  If physics/math use a
%  macro not defined here, add it here (this is the single shared source).
%  files own \documentclass (verified: both real files issue
%  \documentclass[11pt]{article} on line 1), so this preamble must NOT.
% =====================================================================

% ---- encoding / fonts / layout ----
\usepackage[T1]{fontenc}
\usepackage[utf8]{inputenc}
\usepackage[margin=1.1in]{geometry}

% ---- math ----
\usepackage{amsmath,amsthm,amssymb,mathtools}
\usepackage{bm}

% ---- graphics / tables / misc ----
\usepackage{graphicx}
\usepackage{booktabs}
\usepackage{enumitem}
\usepackage{xcolor}

% ---- hyperref last (before xr in the document) ----
\usepackage[colorlinks=true,linkcolor=blue,citecolor=blue,urlcolor=blue]{hyperref}

% =====================================================================
%  Theorem environments — ONE shared numbering scheme across both parts.
%  All share the counter "thmcounter" numbered within section, so Part I
%  and Part II do not collide when cross-referenced.
% =====================================================================
\newtheorem{theorem}{Theorem}[section]
\newtheorem{proposition}[theorem]{Proposition}
\newtheorem{lemma}[theorem]{Lemma}
\newtheorem{corollary}[theorem]{Corollary}
\theoremstyle{definition}
\newtheorem{definition}[theorem]{Definition}
\newtheorem{example}[theorem]{Example}
\theoremstyle{remark}
\newtheorem{remark}[theorem]{Remark}
\newtheorem{note}[theorem]{Note}
\newtheorem{principle}[theorem]{Principle}

% =====================================================================
%  Status labels (README §7: TeX holds only established statements;
%  these tag the epistemic status inline).
% =====================================================================
\newcommand{\established}{\textcolor{green!55!black}{\textsf{[established]}}}
\newcommand{\conditional}{\textcolor{orange!85!black}{\textsf{[conditional]}}}
\newcommand{\deferred}{\textcolor{blue!70!black}{\textsf{[deferred]}}}
\newcommand{\TODO}{\textcolor{red!80!black}{\textsf{[TODO]}}}

% =====================================================================
%  Nonstandard-analysis symbols.
%  *R = hyperreals; st = standard-part map; ord = infinitesimal order;
%  fin = finite hyperreals. eps is the fixed positive infinitesimal
%  lattice spacing (with eps^2 > 0 still infinitesimal).
% =====================================================================
\newcommand{\stR}{{}^{*}\mathbb{R}}        % hyperreals
\newcommand{\stN}{{}^{*}\mathbb{N}}        % hypernaturals
\DeclareMathOperator{\st}{st}              % standard-part map
\DeclareMathOperator{\ord}{ord}            % infinitesimal order
\DeclareMathOperator{\fin}{fin}            % finite hyperreals
\newcommand{\eps}{\varepsilon}

% measurement map  M(f) = st( f / eps^{ord f} )
\newcommand{\Meas}{\mathcal{M}}

% =====================================================================
%  Categorical chain  C_corr -> C_dist -> C_act -> C_eq.
% =====================================================================
\newcommand{\Ccorr}{\mathcal{C}_{\mathrm{corr}}}
\newcommand{\Cdist}{\mathcal{C}_{\mathrm{dist}}}
\newcommand{\Cact}{\mathcal{C}_{\mathrm{act}}}
\newcommand{\Ceq}{\mathcal{C}_{\mathrm{eq}}}

% =====================================================================
%  Common shorthands used across the two parts.
% =====================================================================
\newcommand{\R}{\mathbb{R}}
\newcommand{\N}{\mathbb{N}}
\newcommand{\Z}{\mathbb{Z}}
\newcommand{\C}{\mathbb{C}}
\newcommand{\Q}{\mathbb{Q}}
\newcommand{\corr}{C}                      % correlation C(theta) = -cos theta
\newcommand{\dist}{D}                      % relational distance D = -log|C|
\newcommand{\elliptope}{\mathcal{E}}
\newcommand{\Gram}{G}
\newcommand{\inner}[2]{\langle #1, #2 \rangle}
\newcommand{\abs}[1]{\left| #1 \right|}
\newcommand{\norm}[1]{\left\lVert #1 \right\rVert}
\newcommand{\set}[1]{\left\{ #1 \right\}}
\DeclareMathOperator{\rank}{rank}
\DeclareMathOperator{\tr}{tr}
\DeclareMathOperator{\diag}{diag}
\DeclareMathOperator*{\argmin}{arg\,min}
\DeclareMathOperator*{\argmax}{arg\,max}

% Bell-framework specifics: half-angle delta, Boole quantities F_eps, the
% window depth m = min_eps F_eps.
\newcommand{\Fbool}{F}
\newcommand{\windowdepth}{m}

% =====================================================================
%  Macros required by the real physics.tex / math.tex (verified by macro
%  extraction). Meanings inferred from usage context.
% =====================================================================
\newcommand{\RR}{\mathbb{R}}               % reals (used as \RR)
\newcommand{\ZZ}{\mathbb{Z}}               % integers (Z = pi_1(S^1))
\newcommand{\Cmat}{C}                       % correlation matrix / value, used as |\Cmat|
\newcommand{\dent}{d}                       % relational distance d=-log|C|
\newcommand{\Tcirc}{\mathbb{T}}            % torus; used as \Tcirc^{3}, \Tcirc^{p}
\newcommand{\Scirc}{\mathbb{S}}            % circle S^1; Z \cong pi_1(\Scirc)
\newcommand{\LamQCD}{\Lambda_{\mathrm{QCD}}}% QCD scale
\newcommand{\Eang}{E}                       % E_ang: correlation E(theta) as a function of angle; \Eang(theta)=cos 2theta/(2N)

% end of preamble.tex
 then
%    \usepackage{xr}\externaldocument[I-]{physics}  (in math.tex)
%    \usepackage{xr}\externaldocument[II-]{math}     (in physics.tex)
%    Compile order: bootstrap each once (pdflatex) to make .aux, then
%    latexmk both twice to converge cross-document refs.
%
%  This preamble is deliberately conservative: it defines every macro the
%  two documents are known to reference, plus the common analysis symbols,
%  so a missing \input does not break the build.  If physics/math use a
%  macro not defined here, add it here (this is the single shared source).
%  files own \documentclass (verified: both real files issue
%  \documentclass[11pt]{article} on line 1), so this preamble must NOT.
% =====================================================================

% ---- encoding / fonts / layout ----
\usepackage[T1]{fontenc}
\usepackage[utf8]{inputenc}
\usepackage[margin=1.1in]{geometry}

% ---- math ----
\usepackage{amsmath,amsthm,amssymb,mathtools}
\usepackage{bm}

% ---- graphics / tables / misc ----
\usepackage{graphicx}
\usepackage{booktabs}
\usepackage{enumitem}
\usepackage{xcolor}

% ---- hyperref last (before xr in the document) ----
\usepackage[colorlinks=true,linkcolor=blue,citecolor=blue,urlcolor=blue]{hyperref}

% =====================================================================
%  Theorem environments — ONE shared numbering scheme across both parts.
%  All share the counter "thmcounter" numbered within section, so Part I
%  and Part II do not collide when cross-referenced.
% =====================================================================
\newtheorem{theorem}{Theorem}[section]
\newtheorem{proposition}[theorem]{Proposition}
\newtheorem{lemma}[theorem]{Lemma}
\newtheorem{corollary}[theorem]{Corollary}
\theoremstyle{definition}
\newtheorem{definition}[theorem]{Definition}
\newtheorem{example}[theorem]{Example}
\theoremstyle{remark}
\newtheorem{remark}[theorem]{Remark}
\newtheorem{note}[theorem]{Note}
\newtheorem{principle}[theorem]{Principle}

% =====================================================================
%  Status labels (README §7: TeX holds only established statements;
%  these tag the epistemic status inline).
% =====================================================================
\newcommand{\established}{\textcolor{green!55!black}{\textsf{[established]}}}
\newcommand{\conditional}{\textcolor{orange!85!black}{\textsf{[conditional]}}}
\newcommand{\deferred}{\textcolor{blue!70!black}{\textsf{[deferred]}}}
\newcommand{\TODO}{\textcolor{red!80!black}{\textsf{[TODO]}}}

% =====================================================================
%  Nonstandard-analysis symbols.
%  *R = hyperreals; st = standard-part map; ord = infinitesimal order;
%  fin = finite hyperreals. eps is the fixed positive infinitesimal
%  lattice spacing (with eps^2 > 0 still infinitesimal).
% =====================================================================
\newcommand{\stR}{{}^{*}\mathbb{R}}        % hyperreals
\newcommand{\stN}{{}^{*}\mathbb{N}}        % hypernaturals
\DeclareMathOperator{\st}{st}              % standard-part map
\DeclareMathOperator{\ord}{ord}            % infinitesimal order
\DeclareMathOperator{\fin}{fin}            % finite hyperreals
\newcommand{\eps}{\varepsilon}

% measurement map  M(f) = st( f / eps^{ord f} )
\newcommand{\Meas}{\mathcal{M}}

% =====================================================================
%  Categorical chain  C_corr -> C_dist -> C_act -> C_eq.
% =====================================================================
\newcommand{\Ccorr}{\mathcal{C}_{\mathrm{corr}}}
\newcommand{\Cdist}{\mathcal{C}_{\mathrm{dist}}}
\newcommand{\Cact}{\mathcal{C}_{\mathrm{act}}}
\newcommand{\Ceq}{\mathcal{C}_{\mathrm{eq}}}

% =====================================================================
%  Common shorthands used across the two parts.
% =====================================================================
\newcommand{\R}{\mathbb{R}}
\newcommand{\N}{\mathbb{N}}
\newcommand{\Z}{\mathbb{Z}}
\newcommand{\C}{\mathbb{C}}
\newcommand{\Q}{\mathbb{Q}}
\newcommand{\corr}{C}                      % correlation C(theta) = -cos theta
\newcommand{\dist}{D}                      % relational distance D = -log|C|
\newcommand{\elliptope}{\mathcal{E}}
\newcommand{\Gram}{G}
\newcommand{\inner}[2]{\langle #1, #2 \rangle}
\newcommand{\abs}[1]{\left| #1 \right|}
\newcommand{\norm}[1]{\left\lVert #1 \right\rVert}
\newcommand{\set}[1]{\left\{ #1 \right\}}
\DeclareMathOperator{\rank}{rank}
\DeclareMathOperator{\tr}{tr}
\DeclareMathOperator{\diag}{diag}
\DeclareMathOperator*{\argmin}{arg\,min}
\DeclareMathOperator*{\argmax}{arg\,max}

% Bell-framework specifics: half-angle delta, Boole quantities F_eps, the
% window depth m = min_eps F_eps.
\newcommand{\Fbool}{F}
\newcommand{\windowdepth}{m}

% =====================================================================
%  Macros required by the real physics.tex / math.tex (verified by macro
%  extraction). Meanings inferred from usage context.
% =====================================================================
\newcommand{\RR}{\mathbb{R}}               % reals (used as \RR)
\newcommand{\ZZ}{\mathbb{Z}}               % integers (Z = pi_1(S^1))
\newcommand{\Cmat}{C}                       % correlation matrix / value, used as |\Cmat|
\newcommand{\dent}{d}                       % relational distance d=-log|C|
\newcommand{\Tcirc}{\mathbb{T}}            % torus; used as \Tcirc^{3}, \Tcirc^{p}
\newcommand{\Scirc}{\mathbb{S}}            % circle S^1; Z \cong pi_1(\Scirc)
\newcommand{\LamQCD}{\Lambda_{\mathrm{QCD}}}% QCD scale
\newcommand{\Eang}{E}                       % E_ang: correlation E(theta) as a function of angle; \Eang(theta)=cos 2theta/(2N)

% end of preamble.tex
 then
%    \usepackage{xr}\externaldocument[I-]{physics}  (in math.tex)
%    \usepackage{xr}\externaldocument[II-]{math}     (in physics.tex)
%    Compile order: bootstrap each once (pdflatex) to make .aux, then
%    latexmk both twice to converge cross-document refs.
%
%  This preamble is deliberately conservative: it defines every macro the
%  two documents are known to reference, plus the common analysis symbols,
%  so a missing \input does not break the build.  If physics/math use a
%  macro not defined here, add it here (this is the single shared source).
%  files own \documentclass (verified: both real files issue
%  \documentclass[11pt]{article} on line 1), so this preamble must NOT.
% =====================================================================

% ---- encoding / fonts / layout ----
\usepackage[T1]{fontenc}
\usepackage[utf8]{inputenc}
\usepackage[margin=1.1in]{geometry}

% ---- math ----
\usepackage{amsmath,amsthm,amssymb,mathtools}
\usepackage{bm}

% ---- graphics / tables / misc ----
\usepackage{graphicx}
\usepackage{booktabs}
\usepackage{enumitem}
\usepackage{xcolor}

% ---- hyperref last (before xr in the document) ----
\usepackage[colorlinks=true,linkcolor=blue,citecolor=blue,urlcolor=blue]{hyperref}

% =====================================================================
%  Theorem environments — ONE shared numbering scheme across both parts.
%  All share the counter "thmcounter" numbered within section, so Part I
%  and Part II do not collide when cross-referenced.
% =====================================================================
\newtheorem{theorem}{Theorem}[section]
\newtheorem{proposition}[theorem]{Proposition}
\newtheorem{lemma}[theorem]{Lemma}
\newtheorem{corollary}[theorem]{Corollary}
\theoremstyle{definition}
\newtheorem{definition}[theorem]{Definition}
\newtheorem{example}[theorem]{Example}
\theoremstyle{remark}
\newtheorem{remark}[theorem]{Remark}
\newtheorem{note}[theorem]{Note}
\newtheorem{principle}[theorem]{Principle}

% =====================================================================
%  Status labels (README §7: TeX holds only established statements;
%  these tag the epistemic status inline).
% =====================================================================
\newcommand{\established}{\textcolor{green!55!black}{\textsf{[established]}}}
\newcommand{\conditional}{\textcolor{orange!85!black}{\textsf{[conditional]}}}
\newcommand{\deferred}{\textcolor{blue!70!black}{\textsf{[deferred]}}}
\newcommand{\TODO}{\textcolor{red!80!black}{\textsf{[TODO]}}}

% =====================================================================
%  Nonstandard-analysis symbols.
%  *R = hyperreals; st = standard-part map; ord = infinitesimal order;
%  fin = finite hyperreals. eps is the fixed positive infinitesimal
%  lattice spacing (with eps^2 > 0 still infinitesimal).
% =====================================================================
\newcommand{\stR}{{}^{*}\mathbb{R}}        % hyperreals
\newcommand{\stN}{{}^{*}\mathbb{N}}        % hypernaturals
\DeclareMathOperator{\st}{st}              % standard-part map
\DeclareMathOperator{\ord}{ord}            % infinitesimal order
\DeclareMathOperator{\fin}{fin}            % finite hyperreals
\newcommand{\eps}{\varepsilon}

% measurement map  M(f) = st( f / eps^{ord f} )
\newcommand{\Meas}{\mathcal{M}}

% =====================================================================
%  Categorical chain  C_corr -> C_dist -> C_act -> C_eq.
% =====================================================================
\newcommand{\Ccorr}{\mathcal{C}_{\mathrm{corr}}}
\newcommand{\Cdist}{\mathcal{C}_{\mathrm{dist}}}
\newcommand{\Cact}{\mathcal{C}_{\mathrm{act}}}
\newcommand{\Ceq}{\mathcal{C}_{\mathrm{eq}}}

% =====================================================================
%  Common shorthands used across the two parts.
% =====================================================================
\newcommand{\R}{\mathbb{R}}
\newcommand{\N}{\mathbb{N}}
\newcommand{\Z}{\mathbb{Z}}
\newcommand{\C}{\mathbb{C}}
\newcommand{\Q}{\mathbb{Q}}
\newcommand{\corr}{C}                      % correlation C(theta) = -cos theta
\newcommand{\dist}{D}                      % relational distance D = -log|C|
\newcommand{\elliptope}{\mathcal{E}}
\newcommand{\Gram}{G}
\newcommand{\inner}[2]{\langle #1, #2 \rangle}
\newcommand{\abs}[1]{\left| #1 \right|}
\newcommand{\norm}[1]{\left\lVert #1 \right\rVert}
\newcommand{\set}[1]{\left\{ #1 \right\}}
\DeclareMathOperator{\rank}{rank}
\DeclareMathOperator{\tr}{tr}
\DeclareMathOperator{\diag}{diag}
\DeclareMathOperator*{\argmin}{arg\,min}
\DeclareMathOperator*{\argmax}{arg\,max}

% Bell-framework specifics: half-angle delta, Boole quantities F_eps, the
% window depth m = min_eps F_eps.
\newcommand{\Fbool}{F}
\newcommand{\windowdepth}{m}

% =====================================================================
%  Macros required by the real physics.tex / math.tex (verified by macro
%  extraction). Meanings inferred from usage context.
% =====================================================================
\newcommand{\RR}{\mathbb{R}}               % reals (used as \RR)
\newcommand{\ZZ}{\mathbb{Z}}               % integers (Z = pi_1(S^1))
\newcommand{\Cmat}{C}                       % correlation matrix / value, used as |\Cmat|
\newcommand{\dent}{d}                       % relational distance d=-log|C|
\newcommand{\Tcirc}{\mathbb{T}}            % torus; used as \Tcirc^{3}, \Tcirc^{p}
\newcommand{\Scirc}{\mathbb{S}}            % circle S^1; Z \cong pi_1(\Scirc)
\newcommand{\LamQCD}{\Lambda_{\mathrm{QCD}}}% QCD scale
\newcommand{\Eang}{E}                       % E_ang: correlation E(theta) as a function of angle; \Eang(theta)=cos 2theta/(2N)

% end of preamble.tex

\usepackage{xr}
\externaldocument[II-]{math}

\title{Bell Correlation as a Foundational Principle\\
\large Part I --- Physics: Emergent Spacetime from Bell Correlations}
\author{T. Yanagisawa}
\date{\today}

\begin{document}
\maketitle

\begin{abstract}
We adopt the experimentally established Bell correlations as a foundational
principle and reconstruct emergent spacetime geometry. Taking
$\Eang(\theta)=\cos 2\theta/(2N)$ as the entry point, we show analytically that
the bare framework (correlation $\to$ distance $\to$ multidimensional scaling)
cannot yield a non-compact or single time direction (the \emph{main no-go}), and
that a single topological additional principle (covering~II) generates a single,
non-compact time. We exhibit the minimal $\RR\times\Tcirc^{3}$ spacetime model and
show that continuous rotational $SO(3)$ symmetry emerges at the metric level in the
continuum limit, while the spatial dimension three remains an external input.
\end{abstract}

\tableofcontents

% ======================================================================
\section{Related work and positioning}\label{sec:related}
The idea that spacetime geometry emerges from quantum entanglement or correlations
is well established: the ``entanglement = geometry'' program of Van Raamsdonk~\cite{vanraamsdonk},
Swingle's tensor networks~\cite{swingle}, the Ryu--Takayanagi proposal, Jacobson's
derivation of Einstein's equations from entanglement equilibrium~\cite{jacobson}, and
the It-from-Qubit program. Our correlation$\to$distance$\to$MDS construction sits in
this lineage but is distinguished by fixing a \emph{specific} correlation
$\cos2\theta$ as first principle and proving an analytic no-go.

Two works are closest and require careful delineation.

\paragraph{Milburn (2025)~\cite{milburn}.} Argues that Bell violations are already
evidence for the non-classical causal structure of flat spacetime ``seen from the
inside,'' and that an embedded agent learns geometry from internal measurements (the
``Ringworld'' toy model learns $T(\theta)$ from head-angle settings). This shares our
inside-out, settings-as-vertices stance, but is qualitative: there is no
distance kernel, no MDS metric, no analytic no-go, no covering principle, and no
standard-part measurement. Conversely, Milburn's single-agent retrocausal reading is
absent here.

\paragraph{Karlsson (2018)~\cite{karlsson}.} Closest of all. Eq.~(10) there gives the
same $\cos2\varphi=\mathrm{Re}\,e^{2i\varphi}$, $\sin2\varphi=\mathrm{Im}\,e^{2i\varphi}$
decomposition and the same ``$\pi$-period $\to$ orthogonal axes at $\varphi\in\{0,\pi/4\}$,''
and Eq.~(11) relates 2d and 3d correlations by period doubling $\pi\to2\pi$, arguing
spin-$1/2$ is the unique $D{=}4$ candidate for emergent spatial orientation.
\emph{We acknowledge that the idea ``$\cos2\theta$ yields spatial structure'' is not new.}
Our contribution is the constructive/quantitative development that Karlsson leaves at
the conceptual stage, made precise in \S\ref{sec:nogo}--\S\ref{sec:dim}:
(i) explicit metric recovery via $\dent=-\log|\Cmat|$ and MDS;
(ii) the analytic no-go (Thm.~\ref{thm:nogo});
(iii) covering~II as a \emph{universal cover} $\Scirc\to\RR$ (period $\pi\to\infty$,
non-compact \emph{time}), which is distinct from Karlsson's finite period doubling
$\pi\to2\pi$ (compact, a 2d/3d \emph{space} duality)---see Rmk.~\ref{rem:karlsson};
(iv) the \emph{negative} determination that dimension three is independent of
$\cos2\theta$ (Thm.~\ref{thm:dim}), opposite in direction to Karlsson's positive
claim; and (v) the standard-part measurement theory (Part~II).

% ======================================================================
\section{Foundational principle and method}\label{sec:principle}
\established

We adopt the experimentally established violation of Bell inequalities as a first
principle. The pioneering measurements of Freedman and Clauser~\cite{freedman-clauser},
the time-varying-analyzer experiments of Aspect and
collaborators~\cite{aspect1982,aspect1982tv}, and the strict-locality experiment of
Weihs, Zeilinger and collaborators~\cite{weihs-zeilinger}---recognised by the 2022
Nobel Prize in Physics---establish that nature exhibits correlations no local
hidden-variable theory can reproduce. Rather than seeking to explain these correlations,
we take them as the empirical bedrock and ask what geometry they generate.

Concretely, we fix the two-setting correlation
\begin{equation}
\Eang(\theta)=\frac{\cos 2\theta}{2N},\qquad N=3,
\label{eq:bell}
\end{equation}
where $\theta$ is the relative angle between two measurement settings. The value $N=3$ is
an \emph{identification} with the colour degree of freedom of QCD, not a derivation; we
never claim to derive it. The functional form $\cos2\theta$ is, in the present (maximal)
working stance, part of the adopted principle; Part~II, \S\ref{II-sec:climb} examines how far
it can itself be climbed from a more minimal principle.

The method rests on three pillars. First, the empirical correlation is converted into a
metric quantity through the entanglement distance of Definition~\ref{def:dist}. Second,
the graph on which we work has \emph{measurement settings} as its vertices and
correlation/distance values as its edge data; it is the support of the correlation
matrix, \emph{not} a pre-given lattice of space points. Spatial geometry is an
\emph{output}, obtained by multidimensional scaling (MDS) of the correlation matrix
(\S\ref{sec:rxt3}), never an input. Third, all measurement is modelled as the
standard-part map $\st$ of nonstandard analysis (Part~II, \S\ref{II-sec:st}); Bell
measurement is one instance of this single operation.

\begin{principle}[Bell correlation]\label{prin:bell}
There exists a measurable two-party correlation given by~\eqref{eq:bell}, grounded in the
experimental violations~\cite{freedman-clauser,aspect1982,aspect1982tv,weihs-zeilinger}.
It is taken as given: neither derived nor modelled.
\end{principle}

\begin{definition}[Entanglement distance]\label{def:dist}
For a correlation value $\Cmat$, the \emph{entanglement distance} is
$\dent=-\log|\Cmat|$. Applied to~\eqref{eq:bell}, this converts the empirical correlation
between two settings into a non-negative metric quantity, diverging precisely where the
correlation vanishes.
\end{definition}

% ======================================================================
\section{The main no-go theorem}\label{sec:nogo}
\established\ (notebooks 10--11; analytic backing notebook 22)

The central negative result is that the bare framework of \S\ref{sec:principle}---turning
the correlation~\eqref{eq:bell} into a distance and applying MDS---cannot by itself
produce a non-compact direction or a single time. We isolate three premises, each an
analytic consequence of $\cos2\theta$, and show they jointly forbid both.

\begin{lemma}[Rank-two kernel]\label{lem:rank2}
The kernel $\Cmat(\theta_i,\theta_j)=\cos 2(\theta_i-\theta_j)/(2N)$ factorises, by the
addition formula, as
\begin{equation}
\Cmat(\theta_i,\theta_j)=\sum_{k=1}^{2}\varphi_k(\theta_i)\,\varphi_k(\theta_j),
\qquad \varphi=\frac{1}{\sqrt{2N}}\,(\cos2\theta,\;\sin2\theta).
\label{eq:factor}
\end{equation}
Hence the integral kernel has rank exactly two.
\end{lemma}

\begin{proof}
Expand $\cos2(\theta_i-\theta_j)=\cos2\theta_i\cos2\theta_j+\sin2\theta_i\sin2\theta_j$
and divide by $2N$; this is~\eqref{eq:factor}. The two functions $\cos2\theta,\sin2\theta$
are linearly independent, so the rank is exactly two, not merely at most two.
\end{proof}

\begin{lemma}[Constant norm $\Rightarrow$ bounded image]\label{lem:bounded}
The feature map $u(\theta)=\varphi(\theta)$ of~\eqref{eq:factor} satisfies
$|u(\theta)|^{2}=1/(2N)$, independent of $\theta$; its image lies on a circle of radius
$1/\sqrt{2N}$ and is therefore bounded and compact.
\end{lemma}

\begin{proof}
$|u(\theta)|^2=(\cos^2 2\theta+\sin^2 2\theta)/(2N)=1/(2N)$. We emphasise the logical
separation: Lemma~\ref{lem:rank2} gives only that the image lies in a two-dimensional
subspace (a priori an arbitrary planar curve); the stronger conclusion ``circle'' follows
from the constant-norm identity $\cos^2+\sin^2=1$, a property specific to the form
$\cos2\theta$, not a consequence of rank two alone. For the no-go, boundedness is what is
needed, and boundedness follows from constant norm regardless of the precise shape.
\end{proof}

\begin{lemma}[Period $\pi$ $\Rightarrow$ univalent closed curve]\label{lem:period}
The feature map is $\pi$-periodic, $u(\theta+\pi)=u(\theta)$, so its image is a univalent
closed curve, topologically $\Scirc$.
\end{lemma}

\begin{proof}
$\cos2(\theta+\pi)=\cos(2\theta+2\pi)=\cos2\theta$, and likewise for $\sin$. Thus
$u$ factors through the quotient $\RR/\pi\ZZ\cong\Scirc$, tracing a closed loop once.
\end{proof}

\begin{theorem}[Main no-go]\label{thm:nogo}
Within the bare framework (correlation $\to$ distance $\to$ MDS) applied to~\eqref{eq:bell},
the emergent geometry admits neither a non-compact direction nor a global single time
direction.
\end{theorem}

\begin{proof}
By Lemmas~\ref{lem:rank2} and~\ref{lem:bounded} (premises P1, P2) the MDS image is
contained in the bounded circle of radius $1/\sqrt{2N}$; a bounded image cannot host a
non-compact (unbounded, $\RR$-like) direction. By Lemma~\ref{lem:period} (premise P3) the
image is a univalent closed curve $\Scirc$; assigning an orientation to a closed loop does
not yield a totally ordered, unbounded time axis, and the periodicity returns the curve to
its start, so no global single time direction exists. Both desired features are therefore
excluded. \qedhere
\end{proof}

The three premises P1 (rank two), P2 (constant norm), P3 (period $\pi$) each follow
analytically from $\cos2\theta$; the numerical confirmations of notebooks~10--11 are thus
underwritten by closed-form algebra (notebook~22).

% ======================================================================
\section{Covering principle II: emergence of time}\label{sec:covering}
\established\ (notebook 13)

The no-go theorem shows an additional principle is unavoidable. We adopt a single
topological one, which breaks exactly the premise P3 of Theorem~\ref{thm:nogo} while
preserving P1 and P2.

\begin{principle}[Covering II]\label{prin:cover}
Take the support of the setting space to be not the circle $\Scirc$ but its universal
cover $\RR$, and identify the covering coordinate $t$ (with $\pi:\RR\to\Scirc$,
$t\mapsto t \bmod \pi$) with time.
\end{principle}

\begin{proposition}[Single non-compact time]\label{prop:time}
Under Covering~II the time direction is one-dimensional, totally ordered, and non-compact.
\end{proposition}

\begin{proof}
The universal cover of $\Scirc$ is $\RR$, unique up to isomorphism, with deck
transformation group $\ZZ$. The covering coordinate $t$ ranges over all of $\RR$: it is
totally ordered and unbounded, hence non-compact, and constitutes a single
($1$-dimensional) direction. No continuous parameter is introduced---the deck group is
discrete and the fundamental domain has width fixed by the period of~\eqref{eq:bell}---so
the construction does not smuggle in a tunable infinitesimal (it avoids the ``$\varepsilon$
trap'' discussed in notebook~13).
\end{proof}

\begin{proposition}[Causal transitivity]\label{prop:causal}
The order relation $t_i<t_j$ on the covering coordinate is transitive, yielding a global
causal order.
\end{proposition}

\begin{proof}
The relation is the standard total order on $\RR$, which is transitive. This resolves,
at the structural level, the failure of transitivity encountered when one attempts to
read causal order from the sign of $\cos2\theta$ on the circle (notebook~9, route L-A),
where periodic sign reversals destroy any global order.
\end{proof}

\begin{remark}[Covering~II breaks only P3]\label{rem:p3}
Covering~II leaves the rank-two structure (P1) and the bounded, constant-norm spatial
part (P2) intact: the spatial sector is still the MDS of $|\Cmat|$, an $\Scirc$. Only the
univalence/periodicity premise P3 is removed, and only in the lifted (time) direction.
Thus the operation is the minimal departure from the bare framework consistent with the
no-go analysis (notebook~22, \S22.6).
\end{remark}

\begin{remark}[Covering~II vs.\ Karlsson's period doubling]\label{rem:karlsson}
Covering~II ($\Scirc\to\RR$, period $\pi\to\infty$, non-compact) must not be conflated
with Karlsson's period doubling~\cite{karlsson} ($\pi\to2\pi$, compact). Both alter the
period, but the targets differ decisively: Karlsson's is a finite doubling internal to
\emph{space} (a 2d/3d correlation duality, remaining a compact $\Scirc$), whereas
covering~II is the universal cover producing non-compact \emph{time}. The surface
resemblance (``changing the period'') hides a difference in topology, role, and
output (notebook~23, \S23.3).
\end{remark}

% ======================================================================
\section{The minimal spacetime model $\RR\times\Tcirc^{3}$}\label{sec:rxt3}
\established\ (notebooks 16, 18)

Combining Covering~II with $p$ independent angular settings gives a minimal spacetime.
With $p=4$ (three spatial settings plus the lifted time direction; see \S\ref{sec:dim} for
the status of $p=4$) the model is $\RR\times\Tcirc^{3}$.

\begin{proposition}[Separate construction of time and space]\label{prop:sep}
Time and space must be built separately: the time direction is the covering coordinate
(timelike, dimension exactly one), and the spatial geometry is the MDS of $|\Cmat|$ over
the three angular settings, yielding $\Tcirc^{3}$ of intrinsic dimension three.
\end{proposition}

\begin{proof}
Attempting to embed a single mixed squared interval
$ds^2=-dt^2+\sum d_{\text{space}}^2$ at once produces, via the negative-definite part, a
spurious multiplicity of timelike directions (notebook~16, \S16.1). Constructing time
directly from the covering coordinate makes the timelike dimension exactly one by
construction, while the spatial distances, decomposing additively over the three settings
(Part~II, Prop.~\ref{II-prop:additive}), give an MDS image whose intrinsic dimension is three
(embedding dimension six: three circles, two coordinates each).
\end{proof}

\begin{proposition}[Metric signature and light cone]\label{prop:signature}
The model carries signature $(-,+,+,+)$, with a non-trivial light-cone structure
$ds^2=-dt^2+d_{\text{space}}^2$ in which constant-time slices are everywhere spacelike.
\end{proposition}

\begin{proof}
The covering coordinate contributes one negative (timelike) eigenvalue; the spatial MDS
contributes positive (spacelike) eigenvalues. Sampling pairs of points yields both
timelike ($ds^2<0$) and spacelike ($ds^2>0$) separations, so the light cone is
non-degenerate; at fixed $t$ one has $ds^2=d_{\text{space}}^2\ge0$, so simultaneity slices
are spacelike (notebook~18, part~b). Causal order along $t$ is transitive by
Proposition~\ref{prop:causal}.
\end{proof}

% ======================================================================
\section{Emergence of continuous $SO(3)$ symmetry}\label{sec:so3}
\established\ (notebooks 19, 20)

The discrete construction of $\Tcirc^{3}$ has, as exact isometries, only translations and
the discrete permutation of the three axes; the full continuous rotation group is not an
isometry of the discrete model. We show that continuous $SO(3)$ symmetry nonetheless
emerges at the metric level in the continuum limit, with anisotropy relegated to higher
order.

\begin{theorem}[Metric-level rotational invariance]\label{thm:so3}
Expanding the additive distance kernel for small angular separation $\Delta\varphi$ gives
\begin{equation}
\sum_{c=1}^{3}\dent(\Delta\varphi_c)^2
= \text{const} + 2\,|\Delta\varphi|^2
+ \tfrac{4}{3}\sum_{c}\Delta\varphi_c^{4} + O(|\Delta\varphi|^{6}),
\label{eq:expansion}
\end{equation}
whose leading non-constant term $2|\Delta\varphi|^2$ is spherically symmetric. Hence the
emergent metric is continuously $SO(3)$ invariant, and rotational symmetry breaking first
appears at order $O(|\Delta\varphi|^4)$ as a cubic ($\sum_c\Delta\varphi_c^4$) anisotropy.
\end{theorem}

\begin{proof}
The single-axis kernel admits the small-angle expansion
$-\log(|\cos2\varphi|/(2N)) = \log(2N) + 2\varphi^2 + \tfrac{4}{3}\varphi^4 + O(\varphi^6)$.
Summing over the three axes gives~\eqref{eq:expansion}. The quadratic term is
$2\sum_c\Delta\varphi_c^2=2|\Delta\varphi|^2$, manifestly invariant under $SO(3)$; the
quartic term $\sum_c\Delta\varphi_c^4$ is invariant only under the cubic (octahedral)
subgroup. A control experiment confirms that a genuinely spherical kernel shows no
breaking at any order, isolating the breaking as a property of the additive
(axis-separable) kernel rather than of the lattice discretisation (notebooks~19--20).
\end{proof}

\begin{proposition}[Isotropy originates in the given correlation]\label{prop:isotropy}
The metric-level isotropy of Theorem~\ref{thm:so3} is not circular: it follows from the
homogeneity of the given correlation~\eqref{eq:bell}, not from a cubic lattice put in by
hand.
\end{proposition}

\begin{proof}
Weighting the axes as $\sum_c w_c\,\dent(\Delta\varphi_c)^2$ gives a metric tensor
$\operatorname{diag}(2w_c)$, isotropic if and only if the weights are equal. Equal weights
hold precisely because the same correlation~\eqref{eq:bell} governs every measurement
direction (homogeneity of the given); were the correlation direction-dependent, the metric
would be anisotropic---so the statement is falsifiable, not tautological. The orthogonality
of the axes is a coordinate (gauge) choice with no effect on invariants, while the
\emph{number} of axes (three) is an external input, addressed in \S\ref{sec:dim}
(notebook~20, \S20.6, three-layer separation).
\end{proof}

% ======================================================================
\section{The spatial dimension is an external input}\label{sec:dim}
\established\ (notebooks 15, 17, 21)

The one ingredient the framework does not supply is the spatial dimension itself.

\begin{theorem}[Dimension three is not derivable]\label{thm:dim}
The number of independent settings $p=4$ (equivalently, spatial dimension three) is not
determined within the given correlation~\eqref{eq:bell} together with the standard-part
measurement theory: no condition internal to the framework singles out $p=4$.
\end{theorem}

\begin{proof}
The correlation~\eqref{eq:bell} carries the information of a single relative angle
(rank-two, Lemma~\ref{lem:rank2}); the number of independent settings $p$ is the choice of
how many circles to take in direct product, which does not appear in the correlation
function. Each value $p=1,2,3,\dots$ yields an internally consistent torus $\Tcirc^{p}$ of
intrinsic dimension $p$, with none singled out (notebook~15). The standard-part map places
no bound on $p$ either (Part~II, Rmk.~\ref{II-rem:st-limits}). Hence $p=4$ enters only as an
external input---physically, the number of independent measurement axes of real space
(notebook~17).
\end{proof}

\begin{remark}[Dimension is orthogonal to $\cos2\theta$]\label{rem:orthogonal}
In the stage-2 climbing map (Part~II, \S\ref{II-sec:climb}) the dimension appears at no stage
of the ascent to $\cos2\theta$: the latter closes within a single relative-angle scalar.
Dimension three is therefore not \emph{upstream} of $\cos2\theta$ but an \emph{independent}
axis, contrary to a reading in which spin-$1/2$ correlations would fix $D{=}4$
(cf.\ \S\ref{sec:related}); the origin of dimension three remains open (notebook~21,
\S\S21.5--21.6).
\end{remark}

% ======================================================================
\section{Summary and open problems}\label{sec:summary}
We have taken the experimentally established Bell correlation~\eqref{eq:bell} as a first
principle and shown: (i) an analytic no-go (Theorem~\ref{thm:nogo})---the bare framework
yields neither non-compactness nor a single time; (ii) that a single topological
principle, Covering~II, supplies a single non-compact time
(Propositions~\ref{prop:time}--\ref{prop:causal}) by breaking only premise P3; (iii) the
minimal model $\RR\times\Tcirc^{3}$ with signature $(-,+,+,+)$ and a non-trivial light cone
(\S\ref{sec:rxt3}); (iv) emergence of continuous $SO(3)$ symmetry at the metric level, the
isotropy traceable to the homogeneity of the given correlation
(\S\ref{sec:so3}); and (v) that spatial dimension three is an external input
(Theorem~\ref{thm:dim}), independent of $\cos2\theta$.

Three problems remain open. First, whether the quantum-correlation postulate (A2 of
Part~II, \S\ref{II-sec:climb}) can be obtained from a more minimal information-theoretic
principle. Second, the upstream origin of spatial dimension three, now known to be
independent of $\cos2\theta$. Third (Part~II), the mass-gap argument is established only as
positivity and only under hypothesis~(H1).

% ======================================================================
\begin{thebibliography}{9}
\bibitem{freedman-clauser} S. J. Freedman and J. F. Clauser, \emph{Experimental test of local hidden-variable theories}, Phys. Rev. Lett. \textbf{28}, 938 (1972).
\bibitem{aspect1982} A. Aspect, P. Grangier, and G. Roger, \emph{Experimental realization of Einstein--Podolsky--Rosen--Bohm Gedankenexperiment: a new violation of Bell's inequalities}, Phys. Rev. Lett. \textbf{49}, 91 (1982).
\bibitem{aspect1982tv} A. Aspect, J. Dalibard, and G. Roger, \emph{Experimental test of Bell's inequalities using time-varying analyzers}, Phys. Rev. Lett. \textbf{49}, 1804 (1982).
\bibitem{weihs-zeilinger} G. Weihs, T. Jennewein, C. Simon, H. Weinfurter, and A. Zeilinger, \emph{Violation of Bell's inequality under strict Einstein locality conditions}, Phys. Rev. Lett. \textbf{81}, 5039 (1998).
\bibitem{vanraamsdonk} M. Van Raamsdonk, \emph{Building up spacetime with quantum entanglement}, Gen. Rel. Grav. \textbf{42}, 2323 (2010).
\bibitem{swingle} B. Swingle, \emph{Entanglement renormalization and holography}, Phys. Rev. D \textbf{86}, 065007 (2012).
\bibitem{jacobson} T. Jacobson, \emph{Thermodynamics of spacetime: the Einstein equation of state}, Phys. Rev. Lett. \textbf{75}, 1260 (1995).
\bibitem{milburn} G. J. Milburn, \emph{Spacetime events from the inside out}, arXiv:2503.08715 (2025).
\bibitem{karlsson} A. Karlsson, \emph{Space-time emergence from individual interactions}, arXiv:1806.05710 (2018).
\end{thebibliography}

\end{document}
